%%═══════════════════════════════════════════════════════════════════
%% Lecture 11 — Function Minimization
%% 77315 — Computational Physics
%%═══════════════════════════════════════════════════════════════════

\section{Lecture 11 — Function Minimization}\label{sec:lec11}
\hrule\vspace{1.5em}

%%──────────────────────────────────────────────────────────────────
\subsection*{Overview}
Finding minima of scalar functions arises in energy minimisation,
variational principles, and parameter fitting.  We study three
complementary methods: (1)~\emph{bracketing} via the golden-section
search (derivative-free, 1D); (2)~\emph{gradient descent}
(multidimensional, first-order); (3)~\emph{Newton's method for
minimization} (multidimensional, second-order).

%%──────────────────────────────────────────────────────────────────
\subsection{1D Bracketing: Golden-Section Search}\label{subsec:lec11:golden}

\subsubsection*{Setup}
Given a unimodal function $f(x)$, find three points $a<b<c$ such that
$f(b)<f(a)$ and $f(b)<f(c)$ (a bracket).  Repeatedly probe the
\emph{larger} sub-interval to shrink the bracket.

\subsubsection*{Golden-section ratio}
The probe point $x$ is placed so that both sub-intervals have the same
shrink ratio regardless of which branch is discarded.  This gives the
golden-section ratio
\begin{equation}
  w = \frac{3-\sqrt{5}}{2} \approx 0.382.
  \label{eq:lec11:golden}
\end{equation}
Each iteration reduces the bracket length by factor $1-w = 0.618$
(the golden ratio $\phi$), requiring no function derivatives.

\subsubsection*{Algorithm}
\begin{enumerate}
  \item Find a bracket $(a,b,c)$ with $f(b)<f(a),\,f(b)<f(c)$.
  \item Identify the larger sub-interval (say $[b,c]$).
  \item Place a probe at $x = b + w(c-b)$; evaluate $f(x)$.
  \item If $f(x)<f(b)$: new bracket is $(b,x,c)$; else $(a,b,x)$.
  \item Repeat until $c-a < \eps$.
\end{enumerate}

\nt{The method assumes a single-well (unimodal) shape.  For multimodal functions a global search (e.g.\ random sampling or simulated annealing) is needed to locate a good starting bracket.}

%%──────────────────────────────────────────────────────────────────
\subsection{Gradient Descent}\label{subsec:lec11:gradient}

For a differentiable function $f:\mathbb{R}^n\to\mathbb{R}$, the
steepest-descent update is
\begin{equation}
  \mathbf{x}_{k+1} = \mathbf{x}_k - \alpha_k\,\nabla f(\mathbf{x}_k),
  \label{eq:lec11:gd}
\end{equation}
where $\alpha_k > 0$ is a step length chosen by a line search along
$-\nabla f$.

\subsubsection*{Convergence}
Near the minimum gradient descent converges \emph{linearly} with rate
depending on the condition number $\kappa = \lambda_{\max}/\lambda_{\min}$
of the Hessian.  For ill-conditioned problems conjugate-gradient or
quasi-Newton methods are preferred.

%%──────────────────────────────────────────────────────────────────
\subsection{Newton's Method for Minimization}\label{subsec:lec11:newton}

At a minimum, $\nabla f = \mathbf{0}$.  Apply Newton–Raphson to this
system.  The Hessian $\Hess_{ij} = \partial^2 f/\partial x_i\partial x_j$
plays the role of the Jacobian:
\begin{equation}
  \mathbf{x}_{k+1} = \mathbf{x}_k - \Hess^{-1}(\mathbf{x}_k)\,\nabla f(\mathbf{x}_k).
  \label{eq:lec11:newton}
\end{equation}
Convergence is \emph{quadratic} near the minimum.

\subsubsection*{Positive-definiteness check}
Newton's method only gives a descent direction when $\Hess$ is
positive definite (all eigenvalues $>0$).  A practical check uses
Gaussian elimination: if a negative pivot is encountered, $\Hess$ is
not positive definite and one must fall back to gradient descent for
that step.

\nt{Far from the minimum the Hessian may not be positive definite.  A common fix is to add a multiple of the identity, $\Hess \leftarrow \Hess + \mu I$, (Levenberg–Marquardt damping) to ensure positive definiteness.}

%%──────────────────────────────────────────────────────────────────
\subsection{Summary of Methods}\label{subsec:lec11:summary}

\begin{center}
\begin{tabular}{llll}
\toprule
Method & Derivatives & Convergence & Cost/step \\
\midrule
Golden section  & none        & linear  & 1 eval \\
Gradient descent & $\nabla f$ & linear  & 1 grad \\
Newton           & $\Hess$    & quadratic & $n^2$ \\
\bottomrule
\end{tabular}
\end{center}
